\subsection{শর্ত থেকে গেট নির্বাচন ও সমতা পরীক্ষা}
\chThreeExample{১}{দুটি অনুমতি ছাড়া যন্ত্র চলবে না}
নিরাপত্তা ঢাকনা বন্ধ হলে $A=1$ এবং চালু বোতাম চাপলে $B=1$। দুই শর্ত একসঙ্গে পূরণ হলে মোটর চলবে।
\chThreeAnswer $Y=AB$, তাই AND। যেকোনো একটি শর্ত পূরণ হলেই চালু করতে OR বসালে নিরাপত্তার শর্ত ভঙ্গ হবে। ধারাবাহিকভাবে যুক্ত দুটি আদর্শ সুইচ এই AND আচরণের মডেল।

\chThreeExample{২}{যেকোনো বিপদে অ্যালার্ম}
ধোঁয়া শনাক্তে $S=1$ বা অতিরিক্ত তাপে $T=1$ হলে সংকেত প্রয়োজন।
\chThreeAnswer $Y=S+T$; OR। সমান্তরাল দুটি আদর্শ সুইচে যেকোনো একটি বন্ধ হলে বাতি জ্বলে। প্রকৃত সেন্সর ও মোটর চালানোর জন্য প্রয়োজনীয় বৈদ্যুতিক driver আলাদা অংশ।

\chThreeExample{৩}{চার বিটের দুটি সংখ্যা সমান কি না}
$A=1010$, $B=1000$।
\chThreeAnswer সংশ্লিষ্ট বিটে XNOR করলে ফল $1,1,0,1$। চারটি ফল AND করলে $1\cdot1\cdot0\cdot1=0$; তাই অসমান। শুধু একটি বিট মিললেই পুরো সংখ্যা সমান বলা যাবে না।

\chThreeExample{৪}{তিন-ইনপুট XOR-এর শর্ত}
$1\oplus1\oplus1=1$; তিনটি ইনপুট সমান হলেও XOR আউটপুট ১। দুই-ইনপুট XOR-এর ``ইনপুট ভিন্ন'' নিয়মটি বহু ইনপুটে সরাসরি চলে না। ধারাবাহিক XOR-এ মোট ১-এর সংখ্যা বিজোড় হলে ফল ১; জোড় হলে ০।

\chThreeExample{৫}{ধারণাভিত্তিক MCQ}
দুই ইনপুটে কেবল 00-তে আউটপুট ১ দেয় কোন গেট? (ক) NAND, (খ) NOR, (গ) XOR, (ঘ) AND।
\chThreeAnswer (খ) NOR। OR-এর 00 সারির ০-কে NOT করলে ১ হয়। NAND কেবল 11-তে শূন্য দেয়।

% Mahbub teaching-type follow-up, 2026-10-14
\chThreeExample{৬}{দুই জায়গা থেকে একটি বাতির অবস্থা বদলানো}
একটি আদর্শ control system-এ দুটি switch-এর অবস্থান $A,B$; অবস্থান ভিন্ন হলে বাতি জ্বলে। যেকোনো একটি switch-এর অবস্থান বদলালে বাতির অবস্থা কী হবে?
\chThreeAnswer $Y=A\oplus B$। 00, 01, 10, 11-এ ফল ০, ১, ১, ০। 00 থেকে শুধু B বদলালে 01: বাতি জ্বলে; 01 থেকে শুধু A বদলালে 11: বাতি নেভে। তাই দুই স্থান থেকেই অবস্থা বদলানো যায়। এখানে ০/১ switch-এর \emph{দুটি অবস্থান}; দুটিকে স্বাধীন চালু/বন্ধ অনুমতি ধরে AND-এর নিয়ম বসানো যাবে না।

\chThreeExample{৭}{সুইচের সংযোগ বদলে শর্ত বদলানো}
দুটি আদর্শ switch প্রথমে series-এ আছে: দুটিই বন্ধ হলে বাতি জ্বলে। যেকোনো একটি বন্ধ হলেই জ্বালাতে কী পরিবর্তন করবে?
\chThreeAnswer বন্ধ=১ ধরে প্রথম শর্ত $Y=AB$। parallel সংযোগে $Y=A+B$ হয়। আগে 01 ও 10-এ output ছিল ০, পরে ১; 00 ও 11 অপরিবর্তিত। তাই শুধু 11 input দিয়ে পরীক্ষা করলে পরিবর্তনটি ধরা পড়বে না; সম্পূর্ণ সত্যক সারণি তুলনা করতে হবে।

\chThreeExample{৮}{দুই-input গেট দিয়ে তিন-input কাজ}
শুধু দুই-input গেট দিয়ে $Y=A+B+C$ এবং $Z=ABC$ গঠন করো।
\chThreeAnswer প্রথম OR-এ $P=A+B$; দ্বিতীয় OR-এ $Y=P+C$। প্রথম AND-এ $Q=AB$; দ্বিতীয় AND-এ $Z=QC$। উভয় ক্ষেত্রে দুটি গেট লাগে। $A=1,B=0,C=1$-এ $P=1,Y=1,Q=0,Z=0$। এভাবে associative law ব্যবহার করে বেশি input-এর OR/AND কাজ ছোট গেটে করা যায়। কিন্তু NAND-এর ক্ষেত্রে একইভাবে সরাসরি cascade করলে পূরকের কারণে অন্য রাশি হয়: $\overline{\overline{AB}C}$ সাধারণত $\overline{ABC}$ নয়।
